% ID: SEL-D0062-OB-E02
% ESTADO: validada. S062.
\textbf{a)} Al comenzar, el depósito está lleno y $V(0)=\boxed{8\ \mathrm{m}^3}$.

\textbf{b)} Al completar el cuadrado, $V(t)=(t-16)^2/32$. El depósito se vacía cuando $t=\boxed{16\text{ minutos}}$.

\textbf{c)} En el intervalo físico del vaciado, $0\le t\le16$, la gráfica es el tramo decreciente de la parábola. Pasa por $(0,8)$ y $(8,2)$ y llega al vértice $(16,0)$. Si el depósito sigue vacío después, el volumen se mantiene en cero.
\begin{center}
\setlength{\unitlength}{0.9mm}
\begin{picture}(115,65)(-14,-8)
\qbezier(0,0)(46,0)(92,0)
\qbezier(0,0)(0,24)(0,48)
\put(93,0){$t$}
\put(1,49){$V(t)$}
\qbezier(0,40)(40,0)(80,0)
\put(2,40){\scriptsize $(0,8)$}
\put(40,13){\scriptsize $(8,2)$}
\put(75,-6){\scriptsize $(16,0)$}
\end{picture}
\end{center}
\textbf{d)} $V'(t)=-1+t/16$, por lo que $V'(8)=\boxed{-0{,}5\ \mathrm{m}^3/\mathrm{min}}$. En ese instante, el volumen disminuye a razón de medio metro cúbico por minuto.
